\section{Síntesis y ampliación}

Esta clase extrae de la experiencia de fundación y operación de Okta principios transferibles de la identity infrastructure:

\begin{enumerate}
  \item \textbf{La transition necesita un wedge}: el cambio tecnológico abre una ventana; el acceso estrecho por el que el cliente está dispuesto a pagar hoy se encarga del arranque, y la profundidad de la plataforma, del valor a largo plazo;
  \item \textbf{Identity es el control plane}: directory, authentication, authorization, provisioning, session y audit determinan en conjunto el acceso;
  \item \textbf{La front door exige a la vez seguridad y disponibilidad}: availability, latency, correctness, blast radius y safe change deben validarse en conjunto;
  \item \textbf{La convicción debe conectarse con la evidence}: un entorno de baja probabilidad requiere daily progress, buyer signal y una uncertainty honesta, no rechazar los datos;
  \item \textbf{La gobernanza depende de información sin distorsión}: el CEO y el board deben convertir las malas noticias en options, decision rights y follow-through;
  \item \textbf{Un incident debe cerrar el ciclo de aprendizaje}: containment, RCA, remediation, notification y culture change son todos indispensables;
  \item \textbf{Zero Trust establece una confianza revocable}: el contexto, least privilege, step-up y continuous re-evaluation dependen de un lifecycle limpio;
  \item \textbf{El agent usa delegation}: un principal independiente, actor/subject, short-lived scoped token, vault y audit sustituyen a los secretos compartidos;
  \item \textbf{En una adquisición, primero se protege el customer contract}: compartir primitives no significa fusionar de inmediato el runtime, los equipos y los flujos de trabajo de producto;
  \item \textbf{La IA pasa de la demo al workflow}: el category value requiere que cambien en conjunto las capacidades, los procesos, la organización y los controles de seguridad.
\end{enumerate}

\begin{importantbox}{Tarea de Identity Architecture}
Entrega un diseño de una página para un sistema en el que «un AI agent lee tickets en nombre de un empleado y crea borradores de código»: dibuja al human, el agent principal, el authorization server, el ticket system, el code host y el audit; define el flujo OIDC/OAuth/token-exchange, scope, audience, expiry, vault, session revoke y owner offboarding; después escribe las estrategias de degradación ante IdP outage, stale group, token theft y agent compromise. Queda prohibido que el agent herede directamente los token de larga duración del usuario.
\end{importantbox}

\subsection{Lecturas adicionales}

{\small
\begin{itemize}[nosep,leftmargin=*]
  \item \href{https://developer.okta.com/docs/concepts/universal-directory/}{\textbf{Okta Universal Directory}}: identity data model unificado y source integration.
  \item \href{https://help.okta.com/oie/en-us/content/topics/identity-engine/policies/about-policies.htm}{\textbf{Okta Identity Engine policies}}: global session y app sign-in assurance.
  \item \href{https://sec.okta.com/articles/2023/11/unauthorized-access-oktas-support-case-management-system-root-cause/}{\textbf{Okta support-system incident RCA}}: timeline, logging gap y remediation.
  \item \href{https://csrc.nist.gov/pubs/sp/800/207/final}{\textbf{NIST SP 800-207}}: Zero Trust Architecture.
  \item \href{https://www.rfc-editor.org/rfc/rfc8693}{\textbf{RFC 8693}}: OAuth 2.0 Token Exchange.
  \item \href{https://www.rfc-editor.org/rfc/rfc7644}{\textbf{RFC 7644}}: SCIM provisioning protocol.
  \item \href{https://www.w3.org/TR/webauthn-3/}{\textbf{WebAuthn Level 3}}: phishing-resistant public-key authentication.
\end{itemize}
}

\end{document}
